\documentclass[12pt,a4paper]{article}

% ── Packages ──────────────────────────────────────────────────────────────────
\usepackage[a4paper, margin=2.5cm, headheight=15pt]{geometry}
\usepackage{graphicx}
\usepackage{booktabs}
\usepackage{amsmath}
\usepackage{hyperref}
\usepackage{caption}
\usepackage{subcaption}
\usepackage{float}
\usepackage{parskip}
\usepackage{xcolor}
\usepackage{listings}
\usepackage{fancyhdr}
\usepackage{titlesec}
\usepackage{array}
\usepackage{longtable}
\usepackage{setspace}

% ── Colours ───────────────────────────────────────────────────────────────────
\definecolor{sqlblue}{RGB}{0,70,127}
\definecolor{codegray}{RGB}{245,245,245}
\definecolor{commentgreen}{RGB}{0,128,0}
\definecolor{keywordblue}{RGB}{0,0,200}

% ── SQL Listings Style ────────────────────────────────────────────────────────
\lstdefinestyle{sqlstyle}{
    backgroundcolor=\color{codegray},
    commentstyle=\color{commentgreen}\itshape,
    keywordstyle=\color{keywordblue}\bfseries,
    numberstyle=\tiny\color{gray},
    basicstyle=\ttfamily\small,
    breakatwhitespace=false,
    breaklines=true,
    captionpos=b,
    keepspaces=true,
    numbers=left,
    numbersep=5pt,
    showspaces=false,
    showstringspaces=false,
    showtabs=false,
    tabsize=2,
    frame=single,
    rulecolor=\color{gray!40},
    language=SQL,
    morekeywords={CREATE, OR, REPLACE, VIEW, WITH, OVER, PARTITION, EXTRACT,
                  COALESCE, NULLIF, DATE_TRUNC, ROUND, ROW_NUMBER, RANK, LAG,
                  NUMERIC, TIMESTAMP, FOREIGN, PRIMARY, REFERENCES}
}
\lstset{style=sqlstyle}

% ── Header / Footer ───────────────────────────────────────────────────────────
\pagestyle{fancy}
\fancyhf{}
\rhead{\small Brazilian E-Commerce Analytics}
\lhead{\small Mohammed S. Rahman}
\rfoot{\small \thepage}
\lfoot{\small June 2026}
\renewcommand{\headrulewidth}{0.4pt}
\renewcommand{\footrulewidth}{0.4pt}

% ── Section Formatting ────────────────────────────────────────────────────────
\titleformat{\section}{\large\bfseries\color{sqlblue}}{{\thesection}}{1em}{}[\titlerule]
\titleformat{\subsection}{\normalsize\bfseries}{{\thesubsection}}{1em}{}
\titleformat{\subsubsection}{\normalsize\itshape\bfseries}{{\thesubsubsection}}{1em}{}

% ── Hyperlinks ────────────────────────────────────────────────────────────────
\hypersetup{
    colorlinks=true,
    linkcolor=sqlblue,
    urlcolor=sqlblue,
    citecolor=sqlblue
}

% ── Document ──────────────────────────────────────────────────────────────────
\begin{document}

% ── Title Page ────────────────────────────────────────────────────────────────
\begin{titlepage}
    \centering
    \vspace*{2cm}

    {\Huge\bfseries Brazilian E-Commerce Analytics}\\[0.5em]
    {\Large A Relational Database \& Business Intelligence Pipeline}\\[3cm]

    {\large Mohammed S. Rahman}\\[0.5cm]
    {\normalsize June 2026}\\[4cm]

    \begin{abstract}
        \noindent
        This report documents the end-to-end design, implementation, and visualisation of
        a relational analytics pipeline built on the Olist Brazilian E-Commerce public
        dataset. The pipeline spans database schema creation, data quality validation, the
        construction of reusable analytical views, and a suite of five Power~BI dashboard
        pages covering executive KPIs, revenue trends, product performance, customer
        retention, and delivery \& review quality. Across the dataset's approximately
        100{,}000 orders placed between 2016 and 2018, total platform revenue stands at
        \$15.42~million from 96{,}000 delivered orders, with an average order value of
        \$159.83. Health \& beauty and watches \& gifts are the highest-revenue product
        categories, São Paulo sellers generate the dominant share of platform revenue, and
        the overall late-delivery rate of 8.11\% is associated with a measurable decline in
        customer review scores. The repeat-customer rate of just 3\% highlights a structural
        retention challenge that represents the platform's most significant operational
        opportunity for improvement.
    \end{abstract}

    \vfill
\end{titlepage}

% ── Table of Contents ─────────────────────────────────────────────────────────
\tableofcontents
\newpage

% ══════════════════════════════════════════════════════════════════════════════
\section{Introduction}
% ══════════════════════════════════════════════════════════════════════════════

E-commerce platforms generate rich transactional records that, when organised and queried
systematically, yield actionable insights across revenue, operations, product assortment, and
customer experience. This project applies a structured analytics workflow to the Olist
Brazilian E-Commerce dataset — a publicly available collection of anonymised marketplace
orders — with the goal of building a reproducible, query-driven intelligence layer and a
multi-page Power~BI dashboard accessible to non-technical stakeholders.

The project addresses four core business questions:

\begin{enumerate}
    \item \textbf{Revenue \& Growth:} How has platform revenue evolved over time, which
          states drive the most orders, and what is the trajectory of average order value?
    \item \textbf{Product Performance:} Which product categories generate the most revenue
          and volume, and where does freight cost create margin pressure?
    \item \textbf{Customer Retention:} What share of customers make repeat purchases, and
          how do new-customer acquisition rates evolve across cohorts?
    \item \textbf{Operational Quality:} How does delivery speed correlate with customer
          review scores, and which states and sellers exhibit the worst late-delivery rates?
\end{enumerate}

The remainder of this report is structured as follows. Section~\ref{sec:data} describes the
dataset and the relational schema. Section~\ref{sec:pipeline} documents the SQL pipeline in
detail, from schema creation through analytical views and dashboard views. Section~\ref{sec:dashboard} presents and interprets each of the five Power~BI dashboard pages.
Section~\ref{sec:discussion} synthesises findings and identifies limitations.
Section~\ref{sec:conclusion} concludes.

% ══════════════════════════════════════════════════════════════════════════════
\section{The Data}
\label{sec:data}
% ══════════════════════════════════════════════════════════════════════════════

\subsection{Dataset Overview}

The Olist dataset is a Brazilian e-commerce public dataset made available on Kaggle,
containing real commercial data from the Olist marketplace spanning September 2016 to
August 2018. The dataset captures the full order lifecycle — from customer purchase through
payment, dispatch, delivery, and post-delivery review — across multiple interlinked tables.

Table~\ref{tab:row_counts} summarises the approximate row counts for each table after loading
and quality-checking the data.

\begin{table}[H]
    \centering
    \caption{Approximate row counts per table after data loading.}
    \label{tab:row_counts}
    \begin{tabular}{lrl}
        \toprule
        \textbf{Table} & \textbf{Row Count} & \textbf{Description} \\
        \midrule
        \texttt{customers}                  & 99{,}441 & One record per order-customer link \\
        \texttt{orders}                     & 99{,}441 & Order header with status and timestamps \\
        \texttt{order\_items}               & 112{,}650 & Line items with price and freight \\
        \texttt{order\_payments}            & 103{,}886 & Payment records (multiple per order) \\
        \texttt{order\_reviews}             & 99{,}224 & Post-delivery review scores and comments \\
        \texttt{products}                   & 32{,}951 & Product attributes and dimensions \\
        \texttt{sellers}                    & 3{,}095  & Seller location data \\
        \texttt{product\_category\_translation} & 71 & Portuguese to English category names \\
        \bottomrule
    \end{tabular}
\end{table}

\subsection{Schema Design}

The database uses a star-schema-style relational layout with \texttt{orders} as the central
fact table. Figure~\ref{fig:schema} illustrates the entity relationships.

\begin{figure}[H]
\centering
\ttfamily\small
\begin{tabular}{ccccc}
 & & \textbf{customers} & & \\
 & & $\uparrow$ & & \\
\textbf{order\_reviews} & $\leftarrow$ & \textbf{orders} & $\rightarrow$ & \textbf{order\_payments} \\
 & & $\downarrow$ & & \\
 & & \textbf{order\_items} & & \\
 & $\swarrow$ & & $\searrow$ & \\
\textbf{products} & & & & \textbf{sellers} \\
$\downarrow$ & & & & \\
\textbf{product\_category} & & & & \\
\textbf{\_translation} & & & & \\
\end{tabular}
\rmfamily
\caption{Conceptual entity-relationship diagram. \texttt{orders} sits at the centre,
    joined to \texttt{customers}, \texttt{order\_items}, \texttt{order\_payments}, and
    \texttt{order\_reviews}. \texttt{order\_items} links to both \texttt{products} and
    \texttt{sellers}.}
\label{fig:schema}
\end{figure}

The eight tables and their relationships are defined as follows:

\begin{itemize}
    \item \texttt{customers} — stores one row per \texttt{customer\_id} (which is
          order-scoped in this dataset), with a separate \texttt{customer\_unique\_id} used
          to identify the same buyer across multiple orders.
    \item \texttt{orders} — the primary fact table; each row is a single order with status
          and four key timestamps (purchase, approval, carrier handover, and customer delivery).
    \item \texttt{order\_items} — a composite-keyed table (\texttt{order\_id},
          \texttt{order\_item\_id}) containing the price and freight value for every line item.
    \item \texttt{order\_payments} — records one or more payment records per order,
          supporting split and instalment payments.
    \item \texttt{order\_reviews} — one review per order, containing a 1--5 star score and
          optional comment fields.
    \item \texttt{products} — physical attributes including weight, dimensions, and
          Portuguese-language category names.
    \item \texttt{sellers} — geographic data (city, state, zip code prefix) for each seller.
    \item \texttt{product\_category\_translation} — a lookup table mapping Portuguese
          category names to English equivalents.
\end{itemize}

\subsection{Data Quality Findings}

Before analysis, a structured set of 15 data-quality checks was executed (SQL file
\texttt{02\_data\_quality\_checks.sql}). Key findings are summarised below.

\subsubsection{Order Status Distribution}

The vast majority of orders (approximately 96.5\%) have a status of \texttt{delivered},
making the delivered-order subset the natural population for revenue and performance
analysis. The remaining statuses include \texttt{shipped}, \texttt{canceled},
\texttt{unavailable}, \texttt{invoiced}, \texttt{processing}, and \texttt{created}.

\subsubsection{Missing Delivery Dates}

A meaningful proportion of orders lack a recorded \texttt{order\_delivered\_customer\_date},
consistent with orders that did not reach the \texttt{delivered} status. All revenue and
delivery performance metrics are therefore computed on the delivered subset only, using the
\texttt{vw\_order\_metrics} view which filters explicitly to \texttt{order\_status =
'delivered'}.

\subsubsection{Review Score Distribution}

The platform's review distribution is heavily right-skewed: 5-star reviews are the single
largest group (approximately 57\% of all reviews), while 1-star reviews constitute roughly
11\%. The overall mean review score is 4.09 out of 5.

\subsubsection{Payment Method Distribution}

Credit card is the dominant payment method, accounting for approximately 74\% of all
payment records, followed by \textit{boleto} (a Brazilian bank slip) at roughly 19\%.
Vouchers and debit cards make up the remainder.

% ══════════════════════════════════════════════════════════════════════════════
\section{The Analytics Pipeline}
\label{sec:pipeline}
% ══════════════════════════════════════════════════════════════════════════════

The analytics pipeline consists of nine sequential SQL scripts. Each script builds on the
outputs of its predecessors, progressing from raw table creation through quality checks,
analytical views, domain-specific analyses, and finally dashboard-ready views consumed
directly by Power~BI. Table~\ref{tab:scripts} summarises each script's role.

\begin{table}[H]
    \centering
    \caption{SQL pipeline script inventory.}
    \label{tab:scripts}
    \begin{tabular}{lp{9.5cm}}
        \toprule
        \textbf{Script} & \textbf{Purpose} \\
        \midrule
        \texttt{01\_create\_tables.sql}           & Drops and recreates all eight base tables with
                                                     foreign key constraints. \\
        \texttt{02\_data\_quality\_checks.sql}    & 15 diagnostic queries: row counts, status distributions,
                                                     date coverage, review distributions, and join integrity checks. \\
        \texttt{03\_analytics\_views.sql}         & Creates three core analytical views reused throughout the pipeline:
                                                     \texttt{vw\_order\_metrics}, \texttt{vw\_category\_revenue},
                                                     and \texttt{vw\_customer\_retention}. \\
        \texttt{04\_revenue\_analysis.sql}        & 12 queries covering total revenue, monthly trends, state-level
                                                     breakdowns, month-over-month growth, and AOV. \\
        \texttt{05\_product\_analysis.sql}        & 12 queries on category revenue, volume, freight burden,
                                                     revenue concentration, and review-weighted performance. \\
        \texttt{06\_customer\_retention\_analysis.sql} & 12 queries on repeat-customer rates, cohort retention,
                                                          first-purchase trends, and customer segmentation. \\
        \texttt{07\_delivery\_review\_analysis.sql}    & 16 queries on delivery speed, late rates, review score
                                                          distributions, and the delivery-to-review relationship. \\
        \texttt{08\_seller\_analysis.sql}         & 14 queries on seller revenue, concentration, review
                                                     performance, late-delivery rates, and state-level seller metrics. \\
        \texttt{09\_dashboard\_views.sql}         & 15 \texttt{CREATE OR REPLACE VIEW} statements producing
                                                     the final aggregated views consumed by Power~BI. \\
        \bottomrule
    \end{tabular}
\end{table}

\subsection{Core Analytical Views (Script 03)}

Three foundational views are created in \texttt{03\_analytics\_views.sql} and referenced
extensively by all downstream scripts.

\subsubsection{\texttt{vw\_order\_metrics}}

This is the primary analytical view, joining \texttt{orders}, \texttt{customers}, and
\texttt{order\_items} for all delivered orders. It computes, per order:

\begin{itemize}
    \item \textbf{Revenue components}: \texttt{product\_revenue}, \texttt{freight\_revenue},
          and \texttt{total\_order\_value}.
    \item \textbf{Delivery performance}: \texttt{delivery\_days} (days from purchase to
          customer delivery) and \texttt{is\_late\_delivery} (a binary flag set to 1 when
          actual delivery exceeds the estimated delivery date).
    \item \textbf{Time dimensions}: \texttt{order\_month} via
          \texttt{DATE\_TRUNC('month', \ldots)} for monthly aggregation.
\end{itemize}

A representative excerpt of the view definition is shown below:

\begin{lstlisting}[caption={Core delivered-order metrics view (excerpt).}, label={lst:vw_order}]
CREATE OR REPLACE VIEW vw_order_metrics AS
SELECT
    o.order_id,
    o.customer_id,
    c.customer_unique_id,
    c.customer_state,
    DATE_TRUNC('month', o.order_purchase_timestamp) AS order_month,
    SUM(oi.price)               AS product_revenue,
    SUM(oi.freight_value)       AS freight_revenue,
    SUM(oi.price + oi.freight_value) AS total_order_value,
    CASE
        WHEN o.order_delivered_customer_date
             > o.order_estimated_delivery_date
        THEN 1 ELSE 0
    END AS is_late_delivery,
    EXTRACT(DAY FROM
        o.order_delivered_customer_date
        - o.order_purchase_timestamp) AS delivery_days
FROM orders o
JOIN customers c ON o.customer_id = c.customer_id
JOIN order_items oi ON o.order_id = oi.order_id
WHERE o.order_status = 'delivered'
GROUP BY o.order_id, o.customer_id, c.customer_unique_id,
         c.customer_state, o.order_purchase_timestamp,
         o.order_delivered_customer_date,
         o.order_estimated_delivery_date;
\end{lstlisting}

\subsubsection{\texttt{vw\_category\_revenue}}

This view joins \texttt{order\_items}, \texttt{products}, and
\texttt{product\_category\_translation} to produce revenue, volume, and average price
metrics per product category. The \texttt{COALESCE} pattern handles any categories that
lack an English translation by falling back to the Portuguese name.

\subsubsection{\texttt{vw\_customer\_retention}}

Using a common table expression (CTE) with \texttt{ROW\_NUMBER()} partitioned by
\texttt{customer\_unique\_id}, this view assigns an order sequence number to each delivered
order per unique customer. The outer query then aggregates to flag each customer as a
repeat buyer (\texttt{is\_repeat\_customer = 1}) if they placed more than one delivered
order.

\subsection{Cohort Retention Analysis (Script 06)}

The most technically involved query in the pipeline is the cohort retention analysis in
\texttt{06\_customer\_retention\_analysis.sql}. It uses a four-CTE chain to compute, for
each monthly cohort, the share of customers still active at each subsequent month since
first purchase. The structure is:

\begin{enumerate}
    \item \texttt{customer\_orders} — one row per (unique customer, order month) pair.
    \item \texttt{cohorts} — the minimum order month for each unique customer, defining
          the cohort assignment.
    \item \texttt{cohort\_activity} — cross-joins customers to their cohort to compute
          \texttt{months\_since\_first\_purchase} using \texttt{EXTRACT(YEAR)} and
          \texttt{EXTRACT(MONTH)} arithmetic.
    \item \texttt{cohort\_counts} and \texttt{cohort\_sizes} — aggregate active customer
          counts, then divide by cohort size at month 0 to yield
          \texttt{retention\_percent}.
\end{enumerate}

\subsection{Delivery--Review Relationship (Script 07)}

Script 07 quantifies how delivery speed and on-time performance correlate with customer
satisfaction. Three analytical approaches are used:

\begin{itemize}
    \item \textbf{Late vs on-time review score}: joining \texttt{vw\_order\_metrics} with
          \texttt{order\_reviews} and grouping by \texttt{is\_late\_delivery}.
    \item \textbf{Delivery speed buckets}: orders are bucketed into five groups
          (0--3 days, 4--7 days, 8--14 days, 15--30 days, 31+ days) and the average
          review score computed for each.
    \item \textbf{Review distribution by delivery status}: a window function computes the
          within-status percentage of each review score (1--5), enabling a side-by-side
          comparison of the full review distribution for late versus on-time orders.
\end{itemize}

\subsection{Seller Revenue Concentration (Script 08)}

The seller analysis identifies a concentration pattern in platform revenue. Using
\texttt{RANK()} over seller-level revenue followed by a \texttt{CASE} bucketing step, the
pipeline quantifies what fraction of total revenue is generated by the top 10, top 11--50,
top 51--100, and all remaining sellers.

\subsection{Dashboard Views (Script 09)}

Fifteen \texttt{CREATE OR REPLACE VIEW} statements consolidate the analytical results into
clean, aggregated views designed for direct Power~BI connection via a PostgreSQL
connector. Each view corresponds to one or more visualisations in the dashboard. Named
consistently with a \texttt{dashboard\_} prefix (e.g.\ \texttt{dashboard\_executive\_kpis},
\texttt{dashboard\_cohort\_retention}, \texttt{dashboard\_seller\_concentration}), these
views pre-compute all aggregations so that Power~BI can load data efficiently without
performing complex joins at query time.

\subsection{Data Loading Utility}

The Python script \texttt{load\_reviews.py} handles the loading of the order reviews CSV
into PostgreSQL. The reviews file required special handling due to its
\texttt{latin1} character encoding (common in Brazilian data exports). The script uses
\texttt{pandas} with \texttt{sqlalchemy} and \texttt{psycopg2}, parsing the two timestamp
columns (\texttt{review\_creation\_date} and \texttt{review\_answer\_timestamp}) with
\texttt{pd.to\_datetime(\ldots, errors=`coerce')} to gracefully handle any malformed date
strings. Data is written to the \texttt{order\_reviews} table in chunks of 1{,}000 rows
using the \texttt{multi} insert method.

% ══════════════════════════════════════════════════════════════════════════════
\section{Power BI Dashboard}
\label{sec:dashboard}
% ══════════════════════════════════════════════════════════════════════════════

The Power~BI dashboard comprises five pages, each connected to the corresponding
\texttt{dashboard\_*} views in PostgreSQL. The pages are: Executive Overview, Revenue
(embedded in the Executive Overview), Product Analysis, Customer Retention, and Delivery
\& Reviews, and Seller Analysis.

% ──────────────────────────────────────────────────────────────────────────────
\subsection{Executive Overview}
% ──────────────────────────────────────────────────────────────────────────────

\begin{figure}[H]
    \centering
    \includegraphics[width=\textwidth]{executive_overview}
    \caption{Executive Overview dashboard page showing top-line KPIs, monthly revenue and
    order trends (2016--2018), and revenue by customer state.}
    \label{fig:exec_overview}
\end{figure}

The Executive Overview page (Figure~\ref{fig:exec_overview}) presents six headline KPIs
alongside monthly trend charts and a revenue-by-state breakdown. Key figures are
summarised in Table~\ref{tab:exec_kpis}.

\begin{table}[H]
    \centering
    \caption{Executive KPI summary from \texttt{dashboard\_executive\_kpis}.}
    \label{tab:exec_kpis}
    \begin{tabular}{lr}
        \toprule
        \textbf{Metric} & \textbf{Value} \\
        \midrule
        Total Revenue              & \$15.42 million \\
        Total Delivered Orders     & 96{,}000 \\
        Average Order Value (AOV)  & \$159.83 \\
        Average Review Score       & 4.09 / 5 \\
        Late Delivery Rate         & 8.11\% \\
        Repeat Customer Rate       & 3.00\% \\
        \bottomrule
    \end{tabular}
\end{table}

\subsubsection{Revenue and Order Volume Trends}

Monthly revenue grew substantially from near-zero in late 2016 to approximately \$8--9
million accumulated by mid-2018. The monthly orders chart exhibits a very similar shape,
with a marked acceleration through 2017 that plateaus into 2018. This parallel growth in
both order count and revenue indicates that the average order value remained relatively
stable over the period, a finding confirmed by the AOV metric of \$159.83.

A notable feature of the revenue trend is the apparent levelling off in early-to-mid 2018.
This could reflect seasonal effects, market saturation in the platform's existing customer
base, or data truncation at the dataset's August 2018 endpoint.

\subsubsection{Revenue by State}

The state-level bar chart reveals high geographic concentration. São Paulo (SP) dominates
by a considerable margin, consistent with its status as Brazil's most populous and
economically active state. Rio de Janeiro, Minas Gerais, and other south/southeast states
follow. Northern and northeastern states contribute minimally to platform revenue, which
may reflect both lower internet penetration and longer delivery times to those regions — a
theme reinforced by the Delivery \& Reviews dashboard.

% ──────────────────────────────────────────────────────────────────────────────
\subsection{Product Analysis}
% ──────────────────────────────────────────────────────────────────────────────

\begin{figure}[H]
    \centering
    \includegraphics[width=\textwidth]{product_analysis}
    \caption{Product Analysis dashboard page showing top categories by revenue and items
    sold, revenue share, and freight cost burden.}
    \label{fig:product}
\end{figure}

The Product Analysis page (Figure~\ref{fig:product}) provides four views of category-level
performance sourced from \texttt{vw\_category\_revenue} and
\texttt{dashboard\_product\_categories}.

\subsubsection{Top Categories by Revenue}

Health \& beauty leads all categories at approximately \$1.3 million in product revenue,
followed closely by watches \& gifts (\$1.2 million) and bed, bath \& table (\$1.0 million).
Sports \& leisure and computers \& accessories round out the top five, each contributing
around \$0.9--1.0 million. The top five categories collectively account for roughly 45\% of
total platform revenue — a concentration level that is significant but not extreme, suggesting
that the Olist marketplace has a moderately diversified category mix.

\subsubsection{Top Categories by Items Sold}

By volume (items sold), bed, bath \& table takes the top position at approximately 11,100
items, followed by health \& beauty (9,700 items) and sports \& leisure (8,600 items). The
divergence between the revenue and volume rankings is economically informative: watches
\& gifts ranks second in revenue but seventh in volume, implying a high average selling
price, while bed, bath \& table ranks first in volume but third in revenue, indicating
lower-priced items sold in large quantities.

\subsubsection{Revenue Share}

The revenue share pie chart confirms category concentration: health \& beauty holds
approximately 11\% of total revenue, watches \& gifts around 10.6\%, and the remaining
categories each contribute less than 10\%. The long tail of small categories is visually
represented by the crowded outer portion of the chart, reinforcing that the revenue
distribution is moderately concentrated rather than winner-takes-all.

\subsubsection{Freight Cost Burden by Category}

The freight-to-revenue percentage chart highlights a structurally important margin risk. The
\texttt{home\_comfort} and \texttt{flowers} categories exhibit freight costs exceeding 50\%
of product revenue — an economically unviable position for sellers or the platform. Bulky
and fragile goods (furniture, large home items) unsurprisingly feature high freight burdens.
Categories such as watches \& gifts and computers show moderate freight ratios in the
20--35\% range, reflecting their higher price points relative to weight and volume.

% ──────────────────────────────────────────────────────────────────────────────
\subsection{Customer Retention}
% ──────────────────────────────────────────────────────────────────────────────

\begin{figure}[H]
    \centering
    \includegraphics[width=\textwidth]{customer_retention}
    \caption{Customer Retention dashboard page showing one-time versus repeat customer
    breakdown, new customer acquisition over time, orders-per-customer distribution,
    and cohort retention table.}
    \label{fig:retention}
\end{figure}

The Customer Retention page (Figure~\ref{fig:retention}) is sourced from
\texttt{vw\_customer\_retention}, \texttt{dashboard\_customer\_retention},
\texttt{dashboard\_orders\_per\_customer}, \texttt{dashboard\_new\_customers\_by\_month},
and \texttt{dashboard\_cohort\_retention}.

\subsubsection{One-Time vs Repeat Customers}

The donut chart reveals a striking finding: 97\% of customers (approximately 91,000) made
only one delivered purchase, while just 3\% (approximately 3,000) returned for a second
or subsequent order. This near-complete absence of repeat purchasing behaviour is the most
commercially significant finding in the dataset. It suggests that the platform is operating
almost entirely as a new-customer acquisition engine rather than a retention platform —
a fundamentally expensive model, given that customer acquisition typically costs
significantly more than retention.

\subsubsection{New Customer Acquisition Trend}

The line chart of new customers by year shows rapid growth: near-zero in 2016, a steep
rise through 2017 to approximately 40,000 new customers, and continued growth into 2018
approaching 50,000. This growth trajectory demonstrates strong top-of-funnel performance,
but the retention data contextualises it: virtually none of those new customers returned.

\subsubsection{Orders per Customer Distribution}

The bar chart of customer count by total orders confirms the dominance of single-order
customers, with the bar at \texttt{total\_orders = 1} dwarfing all others by orders of
magnitude. A small but visible tail extends to 5--10+ orders for a tiny segment of highly
engaged customers, but these represent a negligible fraction of the base.

\subsubsection{Cohort Retention Table}

The cohort table (right side of Figure~\ref{fig:retention}) displays the count of customers
returning in each subsequent month after their first purchase, segmented by cohort year.
The 2016 cohort (300 customers at month 0) shows marginal retention to months 1 and
beyond. The 2017 cohort (1,200 at month 0) retains roughly 100 customers at month 1
(approximately 8\%) and rapidly approaches zero by month 3. The 2018 cohort shows similar
or weaker patterns. These figures are consistent with the headline 3\% repeat-customer rate
and confirm that retention is a persistent structural issue rather than a recent trend.

% ──────────────────────────────────────────────────────────────────────────────
\subsection{Delivery \& Reviews}
% ──────────────────────────────────────────────────────────────────────────────

\begin{figure}[H]
    \centering
    \includegraphics[width=\textwidth]{delivery_and_reviews}
    \caption{Delivery \& Reviews dashboard page showing review score distribution,
    delivery speed vs.\ review score, delivery status impact on satisfaction, and
    late-delivery rate by customer state.}
    \label{fig:delivery}
\end{figure}

The Delivery \& Reviews page (Figure~\ref{fig:delivery}) draws from
\texttt{dashboard\_delivery\_reviews}, \texttt{dashboard\_delivery\_speed\_groups},
\texttt{dashboard\_delivery\_by\_state}, and \texttt{dashboard\_review\_distribution}.

\subsubsection{Review Score Distribution}

The histogram of review scores shows a strongly bimodal distribution: 5-star reviews
dominate at approximately 57,000 reviews, with a secondary mode at 1-star (approximately
11,000 reviews). Scores 2, 3, and 4 are less frequent. This polarisation is a common
pattern in marketplace review systems, where customers are most motivated to leave reviews
at the extremes of experience.

\subsubsection{Delivery Speed and Review Score}

The combined bar-and-line chart (top-right of Figure~\ref{fig:delivery}) plots total orders
by delivery speed group alongside the average review score for each group. Results confirm
a clear positive relationship between faster delivery and higher satisfaction:

\begin{itemize}
    \item Orders delivered in \textbf{0--3 days} attract the highest average review scores,
          approaching 4.4 out of 5.
    \item Orders in the \textbf{4--7 day} and \textbf{8--14 day} buckets show moderate
          scores around 4.1--4.2.
    \item Orders taking \textbf{15--30 days} see scores decline to approximately 3.8.
    \item Orders taking \textbf{31+ days} show the lowest satisfaction, with average scores
          around 2.5--2.8.
\end{itemize}

The volume distribution peaks in the 8--14 day bucket, which is the modal delivery
experience on the platform. Reducing this to 4--7 days for a meaningful fraction of orders
would likely produce a measurable improvement in the platform's average review score.

\subsubsection{Late Delivery vs On-Time Review Scores}

The horizontal bar chart confirms that on-time deliveries receive a substantially higher
average review score than late deliveries. On-time orders average approximately 4.3--4.4
stars, while late orders average roughly 2.7--2.9 stars — a difference of approximately 1.5
rating points. This gap has direct implications for seller reputation and platform trust.

\subsubsection{Late Delivery Rate by State}

The horizontal bar chart of late delivery rates by customer state (bottom-right of
Figure~\ref{fig:delivery}) identifies Alagoas (AL), Maranhão (MA), and Piauí (PI) as the
states with the highest proportions of late deliveries, with rates exceeding 20\% for the
worst performers. These states are predominantly in Brazil's northeast region, consistent
with longer last-mile distances from the São Paulo-concentrated seller base and less
developed logistics infrastructure.

% ──────────────────────────────────────────────────────────────────────────────
\subsection{Seller Analysis}
% ──────────────────────────────────────────────────────────────────────────────

\begin{figure}[H]
    \centering
    \includegraphics[width=\textwidth]{seller_analysis}
    \caption{Seller Analysis dashboard page showing top sellers by revenue, seller revenue
    by state, revenue concentration, and seller review performance scatter plot.}
    \label{fig:seller}
\end{figure}

The Seller Analysis page (Figure~\ref{fig:seller}) draws from
\texttt{dashboard\_seller\_performance}, \texttt{dashboard\_seller\_state}, and
\texttt{dashboard\_seller\_concentration}.

\subsubsection{Top Sellers by Revenue}

The top-ten sellers by revenue chart shows the leading seller generating approximately
\$230,000 in product revenue, with a relatively steep drop-off to sellers 2--10, each
contributing in the \$130,000--\$200,000 range. These are large individual operators by
marketplace standards, likely representing either multi-category merchants or high-volume
niche specialists.

\subsubsection{Seller Revenue by State}

São Paulo (SP) is the dominant seller state by a wide margin, with total seller revenue
approaching \$10 million — more than double the combined contribution of all other states.
Paraná (PR) and Minas Gerais (MG) are the next largest seller states. This geographic
concentration of sellers in the south/southeast mirrors the customer concentration in the
same region, and has implications for delivery times to customers in more distant states.

\subsubsection{Revenue Concentration}

The donut chart of seller revenue concentration reveals that the top 10 sellers account for
approximately 19.9\% of total platform revenue (from the 3,000+ active sellers). The top
11--50 sellers contribute a further 13.3\%, and the top 51--100 add 12.3\%, leaving
approximately 54.5\% of revenue generated by the long tail of smaller sellers. This is a
healthier concentration profile than many marketplaces, where a handful of sellers can
dominate 70--80\% of revenue.

\subsubsection{Review Score vs Revenue Scatter}

The scatter plot of average review score against seller revenue (colour-coded by state)
shows a broadly positive cloud without strong linear structure: most sellers cluster at
review scores between 3.5 and 4.5, with revenue varying widely within that band. A small
number of high-revenue sellers show review scores below 3.0, representing a risk to
platform quality. The colour coding reveals that São Paulo sellers (in purple/blue tones)
populate both the high-revenue and the full range of review score positions.

% ══════════════════════════════════════════════════════════════════════════════
\section{Discussion}
\label{sec:discussion}
% ══════════════════════════════════════════════════════════════════════════════

\subsection{Summary of Findings}

This project surfaces five findings of primary commercial importance.

\textbf{1. Revenue growth is strong but decelerating.} Platform revenue grew steeply
through 2017 and into early 2018 before showing signs of deceleration. The average order
value of \$159.83 has remained stable, indicating that growth was volume-driven rather
than ticket-size-driven. Sustaining revenue growth will require either continued new-customer
acquisition or — more efficiently — improvements to the repeat-purchase rate.

\textbf{2. Customer retention is critically low.} A 3\% repeat-customer rate is the
most important single finding in the dataset. The economic implications are severe:
every dollar of revenue growth requires almost exactly one dollar of customer acquisition
spend, and the platform captures almost no benefit from the compounding loyalty that
characterises healthy marketplaces. Initiatives to increase repeat purchasing — loyalty
programmes, post-purchase re-engagement, subscription-style offers — represent the
highest-leverage investment available.

\textbf{3. Delivery speed is the dominant driver of customer satisfaction.} The 1.5-star
average review score gap between late and on-time deliveries, and the monotone decline in
scores across speed buckets, establishes delivery performance as the most operationally
actionable lever for improving customer experience. With 8.11\% of orders arriving late,
and northeastern states suffering rates above 20\%, there is meaningful headroom for
improvement.

\textbf{4. Product concentration is moderate and healthy.} The top five categories
account for roughly 45\% of revenue, and no single category dominates. This diversification
mitigates category risk and suggests the platform has successfully recruited merchants across
a broad assortment. The freight burden analysis, however, flags specific categories
(home comfort, flowers, and oversized goods) where high shipping costs likely suppress
seller viability or require subsidisation.

\textbf{5. Seller geography creates a structural logistics challenge.} The concentration
of both sellers and customers in São Paulo and adjacent states, combined with the low seller
presence in the north and northeast, creates systematically longer and less reliable delivery
routes to those regions. Recruiting sellers in underserved states or establishing
regional fulfilment infrastructure would be the most direct remedies.

\subsection{Limitations and Future Work}

\textbf{Causal claims are not warranted.} The analyses in this project are descriptive
and associative. The observed correlation between late delivery and low review scores does
not rule out the possibility that both are caused by a third factor (e.g., seller quality).
Causal identification would require experimental variation in delivery conditions.

\textbf{The 3\% repeat-customer rate may be partially an artefact of the dataset.}
The Olist dataset's \texttt{customer\_id} is order-scoped, meaning that a returning customer
may receive a new \texttt{customer\_id} if they place an order under slightly different
account details. The \texttt{customer\_unique\_id} field is designed to mitigate this, but
its completeness and accuracy is not formally documented. The true repeat-customer rate may
be modestly higher than 3\%.

\textbf{The dataset ends in August 2018.} Post-2018 growth, competitive dynamics, and
any operational changes Olist made in response to early data signals are outside the scope
of this analysis.

\textbf{Predictive modelling is a natural extension.} The structured database and
analytical views created here provide a clean foundation for supervised learning
applications — for example, predicting which customers are likely to make a second
purchase (a counterpart to the churn prediction pipeline documented separately), or
predicting delivery delay at the order level to trigger proactive customer communications.

% ══════════════════════════════════════════════════════════════════════════════
\section{Conclusion}
\label{sec:conclusion}
% ══════════════════════════════════════════════════════════════════════════════

This project delivers a complete analytics pipeline from raw CSV data to an interactive
multi-page Power~BI dashboard, grounded in a well-structured PostgreSQL relational schema
and a layered SQL architecture. The pipeline produces reproducible, reusable analytical
views that abstract complexity away from the dashboard layer, making future analytical
queries and dashboard extensions straightforward.

The most commercially significant findings are the near-absence of repeat purchasing
behaviour (3\% repeat-customer rate), the strong link between delivery timeliness and
customer satisfaction, and the geographic concentration of sellers in São Paulo that
creates structural delivery disadvantages for northeastern customers. Together, these
findings point toward a clear strategic agenda: invest in customer retention mechanisms,
reduce late-delivery rates in underserved regions, and monitor high-revenue sellers with
sub-par review performance before reputational damage compounds.

The SQL scripts, Python data-loading utility, and Power~BI dashboard represent a
self-contained, extensible artefact that could be maintained and updated as new monthly
order data becomes available. The dashboard views are designed for direct PostgreSQL
connector refresh in Power~BI, requiring no manual data transformation steps once the
underlying tables are updated.

% ══════════════════════════════════════════════════════════════════════════════
% Appendix
% ══════════════════════════════════════════════════════════════════════════════
\appendix

\section{SQL Script: Table Creation (\texttt{01\_create\_tables.sql})}
\label{app:create_tables}

\begin{lstlisting}[caption={Full table creation script.}, label={lst:create_tables}]
DROP TABLE IF EXISTS order_reviews;
DROP TABLE IF EXISTS order_payments;
DROP TABLE IF EXISTS order_items;
DROP TABLE IF EXISTS orders;
DROP TABLE IF EXISTS products;
DROP TABLE IF EXISTS sellers;
DROP TABLE IF EXISTS customers;
DROP TABLE IF EXISTS product_category_translation;

CREATE TABLE customers (
    customer_id           TEXT PRIMARY KEY,
    customer_unique_id    TEXT,
    customer_zip_code_prefix INTEGER,
    customer_city         TEXT,
    customer_state        TEXT
);

CREATE TABLE sellers (
    seller_id             TEXT PRIMARY KEY,
    seller_zip_code_prefix INTEGER,
    seller_city           TEXT,
    seller_state          TEXT
);

CREATE TABLE products (
    product_id            TEXT PRIMARY KEY,
    product_category_name TEXT,
    product_name_lenght   FLOAT,
    product_description_lenght FLOAT,
    product_photos_qty    FLOAT,
    product_weight_g      FLOAT,
    product_length_cm     FLOAT,
    product_height_cm     FLOAT,
    product_width_cm      FLOAT
);

CREATE TABLE orders (
    order_id              TEXT PRIMARY KEY,
    customer_id           TEXT,
    order_status          TEXT,
    order_purchase_timestamp        TIMESTAMP,
    order_approved_at               TIMESTAMP,
    order_delivered_carrier_date    TIMESTAMP,
    order_delivered_customer_date   TIMESTAMP,
    order_estimated_delivery_date   TIMESTAMP,
    FOREIGN KEY (customer_id) REFERENCES customers(customer_id)
);

CREATE TABLE order_items (
    order_id        TEXT,
    order_item_id   INTEGER,
    product_id      TEXT,
    seller_id       TEXT,
    shipping_limit_date TIMESTAMP,
    price           NUMERIC(10, 2),
    freight_value   NUMERIC(10, 2),
    PRIMARY KEY (order_id, order_item_id),
    FOREIGN KEY (order_id)    REFERENCES orders(order_id),
    FOREIGN KEY (product_id)  REFERENCES products(product_id),
    FOREIGN KEY (seller_id)   REFERENCES sellers(seller_id)
);

CREATE TABLE order_payments (
    order_id            TEXT,
    payment_sequential  INTEGER,
    payment_type        TEXT,
    payment_installments INTEGER,
    payment_value       NUMERIC(10, 2),
    FOREIGN KEY (order_id) REFERENCES orders(order_id)
);

CREATE TABLE order_reviews (
    review_id             TEXT PRIMARY KEY,
    order_id              TEXT,
    review_score          INTEGER,
    review_comment_title  TEXT,
    review_comment_message TEXT,
    review_creation_date  TIMESTAMP,
    review_answer_timestamp TIMESTAMP,
    FOREIGN KEY (order_id) REFERENCES orders(order_id)
);

CREATE TABLE product_category_translation (
    product_category_name         TEXT PRIMARY KEY,
    product_category_name_english TEXT
);
\end{lstlisting}

\section{SQL Script: Customer Retention View (\texttt{03\_analytics\_views.sql} excerpt)}
\label{app:retention_view}

\begin{lstlisting}[caption={Customer retention view using \texttt{ROW\_NUMBER()} to identify repeat purchasers.}, label={lst:retention_view}]
CREATE OR REPLACE VIEW vw_customer_retention AS
WITH customer_orders AS (
    SELECT
        c.customer_unique_id,
        o.order_id,
        o.order_purchase_timestamp,
        ROW_NUMBER() OVER (
            PARTITION BY c.customer_unique_id
            ORDER BY o.order_purchase_timestamp
        ) AS order_number
    FROM orders o
    JOIN customers c ON o.customer_id = c.customer_id
    WHERE o.order_status = 'delivered'
)
SELECT
    customer_unique_id,
    COUNT(order_id)                           AS total_orders,
    MIN(order_purchase_timestamp)             AS first_order_date,
    MAX(order_purchase_timestamp)             AS last_order_date,
    CASE
        WHEN COUNT(order_id) > 1 THEN 1
        ELSE 0
    END AS is_repeat_customer
FROM customer_orders
GROUP BY customer_unique_id;
\end{lstlisting}

\section{Python Data Loading Script (\texttt{load\_reviews.py})}
\label{app:load_reviews}

\begin{lstlisting}[language=Python,
    caption={Python script for loading the order reviews CSV into PostgreSQL, handling Latin-1 encoding and timestamp parsing.},
    label={lst:load_reviews},
    keywordstyle=\color{keywordblue}\bfseries,
    stringstyle=\color{red!60!black},
    commentstyle=\color{commentgreen}\itshape]
import pandas as pd
from sqlalchemy import create_engine
from urllib.parse import quote_plus

DB_USER     = "****"
DB_PASSWORD = "****"
DB_HOST     = "localhost"
DB_PORT     = "****"
DB_NAME     = "ecommerce_analytics"

csv_path = r"C:\Users\..."
encoded_password = quote_plus(DB_PASSWORD)

engine = create_engine(
    f"postgresql+psycopg2://"
    f"{DB_USER}:{encoded_password}@{DB_HOST}:{DB_PORT}/{DB_NAME}"
)

df = pd.read_csv(csv_path, encoding="latin1", dtype=str)

date_cols = ["review_creation_date", "review_answer_timestamp"]
for col in date_cols:
    df[col] = pd.to_datetime(df[col], errors="coerce")

df.to_sql(
    "order_reviews", engine,
    if_exists="append", index=False,
    chunksize=1000, method="multi"
)
print("Loaded order_reviews successfully.")
\end{lstlisting}

\end{document}
